\documentclass[11pt]{article}
\usepackage[margin=1in]{geometry}
\usepackage[T1]{fontenc}
\usepackage{graphicx}
\usepackage{hyperref}
\usepackage{enumitem}
\setlength{\parskip}{0.5em}
\setlength{\parindent}{0pt}

\title{test-time compute and reasoning in large language models}
\author{arXiv GraphRAG --- phase-6 deep research}
\date{\today}

\begin{document}
\maketitle

\begin{abstract}
(judge output unparseable)
\end{abstract}

\section{Scaling the Moment of Thought: A Review of Test-Time Compute and Reasoning in Large Language Models}

For much of the recent history of artificial intelligence, the advancement of large language models (LLMs) has been defined by scaling laws centered on pre-training compute, parameter count, and dataset volume. However, as models encounter increasingly complex logical and mathematical challenges, a fundamental limitation has become apparent: the reliance on "fast," intuitive next-token prediction often fails to support deep, multi-step deliberation. This has necessitated a paradigm shift from scaling only during the training phase to the strategic allocation of computational resources during the generation phase. Understanding test-time compute is critical because it represents a new dimension of intelligence, allowing models to trade increased inference latency for improved reasoning accuracy, effectively moving from reactive pattern matching to proactive problem-solving.

This review traces the evolution of this capability, beginning with the foundational principles found in the Foundations of LLM Scaling and Reasoning to establish how intelligence was traditionally viewed through the lens of training-time scaling. We then examine the transition toward Early Reasoning Paradigms: Prompting and Verification, where initial attempts to elicit reasoning were made through heuristic methods like Chain-of-Thought. The narrative then shifts to the contemporary landscape marked by The Rise of Inference-Time Scaling and o1-like Models, which formalizes the use of extra computation at the moment of inference. Following a Systematic Analysis and Survey of Test-Time Compute, we explore the technical methodologies for Optimizing Test-Time Compute: Efficiency, Search, and Agents, detailing how search algorithms and agentic loops refine the reasoning process. Finally, we look toward the Frontiers: Domain-Specific Scaling and Resource Constraints, discussing how specialized tasks and hardware limitations will define the next era of deliberative AI.

\subsection{Foundations of LLM Scaling and Reasoning}

The evolution of large language models (LLMs) is rooted in the transition from specialized architectures to generalized, scale-driven systems. The foundational shift began with the introduction of the Transformer architecture (\texttt{arXiv:1706.03762}), which replaced recurrent structures with attention mechanisms. This architecture enabled significantly higher resource efficiency through parallelization, such as achieving a BLEU score of 41.8 for English-to-French translation in just 3.5 days on eight GPUs (\texttt{arXiv:1706.03762}). Building upon this, research established that model performance, measured by cross-entropy loss, follows predictable power-law scaling behaviors relative to model size, dataset size, and total compute (\texttt{arXiv:2001.08361}). Crucially, these scaling laws suggest that larger models exhibit superior sample efficiency, often necessitating a strategic balance between model scale and data volume to optimize training (\texttt{arXiv:2001.08361}).

This scaling paradigm was operationalized by GPT-3 (\texttt{arXiv:2005.14165}), which demonstrated that increasing parameters up to 175 billion enables task-agnostic, few-shot learning across diverse domains, including 3-digit arithmetic and news generation. While GPT-3 showed competitiveness with prior fine-tuning methods (\texttt{arXiv:2005.14165}), its performance across massive multitask benchmarks (\texttt{arXiv:2009.03300}) revealed a significant gap: although the largest model improved over random chance by nearly 20 percentage points across 57 tasks—ranging from US history to computer science—it still struggled to reach expert-level accuracy (\texttt{arXiv:2009.03300}).

Subsequent research has focused on optimizing this scaling by refining training data and model efficiency. The LLaMA series (\texttt{arXiv:2302.13971}) demonstrated that models trained on trillions of tokens can achieve superior performance with fewer parameters; for instance, LLaMA-13B outperformed the much larger GPT-3 (175B) on most benchmarks. This trajectory toward efficient, open foundation models continued with Llama 2 (\texttt{arXiv:2307.09288}), which extended these capabilities into dialogue-specific use cases through fine-tuning for helpfulness and safety (\texttt{arXiv:2307.09288}).

While much of this literature focuses on scaling pre-training parameters, a distinct shift toward improving reasoning through inference-time mechanisms is evident. Rather than relying solely on model size, the "Self-Consistency" method (\texttt{arXiv:2203.11171}) introduces a decoding strategy that samples diverse reasoning paths and selects the most consistent answer via marginalization. This approach provides substantial accuracy gains in arithmetic and commonsense reasoning, such as a 17.9\% improvement on GSM8K and 12.2\% on AQuA (\texttt{arXiv:2203.11171}). Collectively, these works illustrate a progression from establishing the architectural efficiency of attention (\texttt{arXiv:1706.03762}) and the predictability of scale (\texttt{arXiv:2001.08361}) to the refinement of both parameter-heavy foundation models (\texttt{arXiv:2302.13971}) and specialized reasoning strategies (\texttt{arXiv:2203.11171}).

\subsection{Early Reasoning Paradigms: Prompting and Verification}

Early approaches to enhancing reasoning in large language models (LLMs) focused on eliciting latent capabilities through specialized prompting techniques and external verification mechanisms. A foundational shift occurred with the introduction of "scratchpads," where models emit intermediate computation steps to facilitate multi-step tasks such as long addition or program execution [\texttt{arXiv:2112.00114}]. This concept was significantly expanded by Chain-of-Thought (CoT) prompting [\texttt{arXiv:2201.11903}], which provides intermediate reasoning exemplars to elicit emergent reasoning abilities in sufficiently large models. CoT has since been refined through several specialized prompting strategies: "Least-to-Most" prompting decomposes complex problems into simpler subproblems to enable easy-to-hard generalization [\texttt{arXiv:2205.10625}], while "Zero-Shot CoT" utilizes the simple phrase "Let's think step by step" to improve performance on arithmetic and logical tasks [\texttt{arXiv:2205.11916}].

Beyond simple prompting, researchers introduced methods to increase test-time compute by exploring multiple reasoning paths. Self-Consistency [\texttt{arXiv:2203.11171}] improves upon CoT by sampling diverse reasoning paths and selecting the most consistent answer via marginalization, yielding significant accuracy gains across datasets like GSM8K (+17.9\%) and SVAMP (+11.0\%). Moving toward more deliberate decision-making, the Tree of Thoughts (ToT) framework [\texttt{arXiv:2305.10601}] enables exploration over coherent "thoughts" using lookahead and backtracking, demonstrating a massive performance leap in tasks like the Game of 24 (74\% for ToT vs. 4\% for CoT). Alternatively, Program-aided Language Models (PAL) [\texttt{arXiv:2211.10435}] offload the computational burden to a symbolic Python interpreter, allowing the LLM to focus on problem decomposition; this synergy allows smaller models to surpass massive architectures like PaLM-540B on GSM8K.

A parallel line of research focuses on verification to mitigate reasoning errors. Early work utilized trained verifiers to rank multiple candidate solutions [\texttt{arXiv:2110.14168}], a method that scales more effectively with data than standard finetuning. This evolved into more granular "step-aware" verification. While some approaches focus on outcome-based feedback [\texttt{arXiv:2211.14275}], newer paradigms prioritize process supervision—providing feedback at the individual step level [\texttt{arXiv:2305.20050}]. For instance, the DIVERSE framework [\texttt{arXiv:2206.02336}] employs a step-by-step verifier, and Math-Shepherd [\texttt{arXiv:2312.08935}] utilizes a Process Reward Model to both rerank outputs and reinforce models via Proximal Policy Optimization (PPO), achieving up to 89.1\% accuracy on GSM8K. Collectively, these paradigms transition from simple prompt-based elicitation to sophisticated test-time search and multi-level verification strategies.

\subsection{The Rise of Inference-Time Scaling and o1-like Models}

Recent advancements in large language models (LLMs) have shifted focus from purely scaling model parameters to optimizing test-time compute to enhance reasoning capabilities. This paradigm, often characterized by "o1-like" behavior, utilizes extended chain-of-thought (CoT) processes and reinforcement learning (RL) to solve complex tasks in mathematics, coding, and STEM fields [\texttt{arXiv:2412.16720}, \texttt{arXiv:2501.12948}].

A primary mechanism for this scaling involves incentivizing emergent reasoning patterns—such as self-reflection and dynamic strategy adaptation—through pure RL [\texttt{arXiv:2501.12948}]. For instance, DeepSeek-R1 demonstrates that RL can drive superior performance on verifiable tasks compared to traditional supervised learning [\texttt{arXiv:2501.12948}]. This aligns with the broader observation that scaling inference compute can be more computationally efficient than scaling model parameters, with smaller models employing advanced inference algorithms offering Pareto-optimal trade-offs [\texttt{arXiv:2408.00724}].

Various strategies have been proposed to implement this scaling. Some approaches utilize AlphaZero-like tree-search frameworks, which employ learned value functions during both training and decoding to guide decision-making [\texttt{arXiv:2309.17179}]. Others rely on repeated sampling, where coverage scales with the number of samples, often following a log-linear relationship [\texttt{arXiv:2407.21787}]. However, simple sampling methods like majority voting and reward-model-based selection tend to plateau after several hundred samples [\texttt{arXiv:2407.21787}]. To overcome such plateaus, more sophisticated search methods are required. For example, searching against dense, process-based verifier reward models can improve efficiency by over 4x compared to best-of-N baselines [\texttt{arXiv:2408.03314}]. Such compute-optimal strategies can even allow test-time compute to outperform models that are 14x larger in FLOPs-matched evaluations [\texttt{arXiv:2408.03314}].

Recent work also explores "budget forcing" as a simple scaling mechanism, where the thinking process is forcefully terminated or extended to extrapolate performance on difficult benchmarks like AIME24 [\texttt{arXiv:2501.19393}]. While these methods enhance reasoning, they introduce the risk of "overthinking." Research indicates that for simple problems (e.g., basic arithmetic), extended CoT processes can lead to excessive computational overhead with minimal benefit, suggesting a need for more rational resource allocation [\texttt{arXiv:2412.21187}]. This tension between scaling for complexity and maintaining efficiency for simplicity remains a central challenge in the development of next-generation reasoning models.

\subsection{Systematic Analysis and Survey of Test-Time Compute}

The landscape of test-time compute (TTC) can be categorized through various structural and operational frameworks. From a structural perspective, scaling can be viewed as the allocation of compute via subproblem decomposition, utilizing sequential, parallel, or tree-structured topologies such as Tree-of-Thought [\texttt{arXiv:2511.14772}]. Operationally, researchers distinguish between L1-controllability, which involves fixed compute budgets, and L2-adaptiveness, which dynamically scales computation based on model confidence or input difficulty [\texttt{arXiv:2507.02076}].

A critical tension emerges regarding the efficacy of scaling: while increasing compute is often assumed to improve performance, recent evidence suggests significant limitations and non-monotonic behaviors. For instance, scaling through longer Chain-of-Thought (CoT) can actually impair reasoning performance in specific domains, suggesting that there exists an optimal scaled length distribution rather than a "more is better" paradigm [\texttt{arXiv:2502.18080}]. This is further complicated by the fact that increasing token counts does not necessarily translate to higher accuracy in highly complex regimes [\texttt{arXiv:2504.00294}]. Comprehensive benchmarks across diverse tasks—including arithmetic, logical, and planning reasoning—confirm that no single scaling technique performs consistently well across all task categories [\texttt{arXiv:2502.12521}]. Furthermore, the relationship between training and inference is nuanced; using standard cross-entropy loss can lead to model overconfidence, which paradoxically decreases pass@N accuracy during longer training phases [\texttt{arXiv:2502.07154}].

To mitigate these inefficiencies, research has shifted toward optimization and adaptive strategies. Rather than employing a fixed number of samples, "Self-Calibration" techniques distill Self-Consistency-derived confidence to enable early-stopping in Best-of-N and Self-Consistency frameworks [\texttt{arXiv:2503.00031}]. This approach improves resource efficiency by adjusting computation based on query complexity; for example, in MathQA tasks, this method improved accuracy from 81.0 to 83.6 using a sample budget of 16 [\texttt{arXiv:2503.00031}]. Similarly, "Thinking-Optimal Scaling" seeks to teach models varying reasoning efforts by using seed data with diverse response length distributions, ultimately selecting the shortest correct response for self-improvement [\texttt{arXiv:2502.18080}]. These advancements suggest that the future of test-time scaling lies not in brute-force computation, but in the precise, adaptive calibration of reasoning effort to the complexity of the task at hand.

\subsection{Optimizing Test-Time Compute: Efficiency, Search, and Agents}

Recent advancements in test-time scaling (TTS) shift the focus from increasing model parameters to optimizing the allocation of computational resources during inference. A central theme in this literature is the pursuit of compute-optimality, where the effectiveness of scaling is highly contingent upon the interplay between the policy model, the Process Reward Model (PRM), and the inherent difficulty of the task [\texttt{arXiv:2502.06703}]. This relationship is so potent that smaller models, such as a 1B parameter LLM, can potentially surpass a 405B parameter model when utilizing an optimal TTS strategy [\texttt{arXiv:2502.06703}]. Furthermore, research indicates that for a given model type, optimal TTS performance scales monotonically with the provided compute budget [\texttt{arXiv:2512.02008}].

Scaling mechanisms vary significantly across the literature, ranging from architectural modifications to complex search strategies. One approach involves latent reasoning, where a single model utilizes a recurrent block to unroll to arbitrary depth in latent space, improving performance up to a computation load equivalent to 50 billion parameters [\texttt{arXiv:2502.05171}]. In contrast, search-based methods focus on structured exploration. For example, hierarchical retrieval-augmented MCTS integrates dual-level retrieval (coarse templates and fine-grained steps) with MCTS and PRMs, yielding a 16\% relative improvement for an 8B model on mathematical reasoning [\texttt{arXiv:2507.05557}]. Similarly, dual-phase search frameworks separate planning from execution, using reward models to prune paths and dynamically reallocate budgets to challenging segments, thereby reducing redundant computation [\texttt{arXiv:2509.25420}].

The literature also highlights a transition toward agentic and multi-agent paradigms to handle complex, multi-stage tasks. While scaling test-time compute can improve general language agent performance through parallel sampling and sequential revision [\texttt{arXiv:2506.12928}], more specialized frameworks like AgentTTS use LLM agents to autonomously search for optimal model and budget allocations through iterative feedback [\texttt{arXiv:2508.00890}]. This agentic approach addresses the impracticality of brute-force search due to high inference costs [\texttt{arXiv:2508.00890}]. Moving further, TMAS utilizes multi-agent synergy—combining specialized agents with hierarchical memories (experience and guideline banks)—to achieve stronger iterative scaling than existing baselines [\texttt{arXiv:2605.10344}].

Despite these gains, efficiency remains a critical bottleneck. Speculative decoding is proposed as a method to mitigate the overhead caused by repetitive reasoning traces in Best-of-N sampling or multi-round thinking, with n-gram-based methods proving particularly effective at capturing repetitive patterns [\texttt{arXiv:2509.04474}]. Ultimately, the field is moving toward model- and task-aware strategies [\texttt{arXiv:2506.13102}], where the goal is to balance the latency and cost-per-token against the reasoning requirements of the specific domain [\texttt{arXiv:2509.19645}].

\subsection{Frontiers: Domain-Specific Scaling and Resource Constraints}

The efficacy of test-time compute (TTC) is not a universal constant; rather, it is highly contingent upon the task domain, the model's underlying knowledge, and the structural method of allocation. In structured reasoning domains such as mathematics, STEM, and NP-hard problems, scaling mechanisms—including lengthening generated scratchpads and utilizing verifiers—provide clear advantages, though these benefits diminish as problem complexity increases (\texttt{arXiv:2504.00294}). Similarly, in technical domains like RTL code generation, scaling both reasoning data and TTC through iterative self-correction can yield significant accuracy improvements (\texttt{arXiv:2506.05566}).

However, a critical tension exists between reasoning depth and task type, often manifesting as "overthinking" or hallucination. In medical reasoning, performance improves with increased token budgets but reaches an optimal limit at approximately 4K tokens, beyond which performance may degrade (\texttt{arXiv:2504.00869}). This risk is even more acute in knowledge-intensive tasks, where increasing compute-only scaling fails to improve accuracy and instead increases hallucinations, as the model cannot generate information beyond what is already encoded in its parameters (\texttt{arXiv:2509.06861}). This suggests a fundamental boundary for TTC: it functions as a mechanism for logical processing rather than information retrieval. This limitation is reinforced by evidence that if the necessary skills are absent from the initial training data, increasing test-time compute can actually harm performance (\texttt{arXiv:2510.03605}).

The transition from general to domain-specific scaling further reveals significant performance plateaus. While general-purpose reasoning models show inconsistent benefits in machine translation, domain-specific fine-tuning allows for consistent improvements up to an optimal reasoning depth (\texttt{arXiv:2510.06471}). These saturation points are observed across various scaling dimensions: increasing the number of samples (\$K\$) in temperature-based sampling eventually yields no further gains (\texttt{arXiv:2510.02611}), and general reasoning benchmarks eventually reach points where additional computation yields diminishing returns (\texttt{arXiv:2505.20522}).

To mitigate these constraints, research is shifting from static scaling toward resource-efficient, structured, and dynamic allocation. Rather than applying uniform compute, dynamic frameworks that select between methods like best-of-N or incremental decoding achieve superior accuracy-cost and latency trade-offs (\texttt{arXiv:2509.09864}). This is supported by a move toward structured subproblem decomposition, utilizing sequential, parallel, or tree-structured topologies (\texttt{arXiv:2511.14772}). Ultimately, the deployment of TTC must balance predictive accuracy against the energy costs of additional computation (\texttt{arXiv:2505.14733}) and the potential for performance degradation in open-source vision-language models during iterative refinement (\texttt{arXiv:2512.11109}).

\subsection{Tensions and Open Problems}

A fundamental debate persists regarding the optimal mechanism for enhancing model performance: whether intelligence is best derived from massive pre-training or through the strategic allocation of computational resources during inference. Early paradigms, such as those established in \texttt{arXiv:2005.14165}, focused on the emergent few-shot capabilities of scaling parameters. However, more recent developments, exemplified by the GPT-4 technical report in \texttt{arXiv:2303.08774}, suggest that the frontier of capability is increasingly defined by how models interact with complex prompts and iterative processes. This tension is further complicated by the emergence of test-time scaling. While some argue that scaling inference compute via repetitive sampling or "large language monkeys" can bridge performance gaps, as discussed in \texttt{arXiv:2407.21787}, others contend that this approach is inefficient. For instance, \texttt{arXiv:2502.06703} argues that compute-optimal testing requires a more nuanced balance, suggesting that a smaller, highly optimized model can potentially surpass a massive 405B parameter model if test-time compute is managed more effectively.

A second tension concerns the reliability and methodology of self-correction and iterative refinement. One school of thought posits that models can autonomously improve their outputs through verbal reinforcement or self-feedback loops, as seen in the "Reflexion" framework of \texttt{arXiv:2303.11366} and the "Self-Refine" approach in \texttt{arXiv:2303.17651}. These methods aim to move beyond the static instruction-following paradigms established in \texttt{arXiv:2203.02155} by allowing models to critique their own reasoning. Conversely, critics suggest that these iterative processes may be fundamentally limited by the model's underlying calibration. This is highlighted in \texttt{arXiv:2503.00031}, which critiques the broad assumptions of the GPT-4 report by arguing for more efficient, self-calibrated scaling methods that do not rely on brute-force iteration. Furthermore, the efficacy of these refinement loops is often contested by the need for structured reasoning; while \texttt{arXiv:2203.11171} advocates for self-consistency to improve Chain of Thought, more complex architectures like the "Tree of Thoughts" in \texttt{arXiv:2305.10601} suggest that simple consistency is insufficient for deliberate, multi-path problem solving.

Finally, there is an unresolved conflict regarding the "scaling laws" of reasoning: specifically, whether increasing test-time compute leads to a permanent gain in intelligence or a diminishing return plateau. While current trends emphasize scaling inference via multi-agent synergy (\texttt{arXiv:2605.10344}) or latent recurrent depth (\texttt{arXiv:2502.05171}), there are growing concerns regarding the sustainability of these methods. As noted in \texttt{arXiv:2505.20522}, researchers are beginning to observe a "scaling plateau" where additional test-time compute yields marginal improvements in reasoning models. This is compounded by the practical constraints of latency and energy consumption, as explored in \texttt{arXiv:2510.02611} and \texttt{arXiv:2505.14733}, which suggest that the path forward may not be through more compute, but through more \textit{intelligent} compute allocation.

\subsubsection{Outlook}

The field is currently transitioning from a "scaling-at-all-costs" era to one of "inference-time efficiency." The frontier is moving toward specialized, task-aware scaling strategies—particularly in high-stakes domains like medical reasoning (\texttt{arXiv:2504.00869})—where the goal is to maximize reasoning depth while minimizing token overhead and energy expenditure. As we move toward 2026, the focus will likely shift from simply increasing the number of samples to developing sophisticated, system-level architectures that can dynamically decide when to "think harder," effectively treating test-time compute as a finite, precious resource to be optimized through latent reasoning and multi-agent coordination.

\section*{References}
\begin{enumerate}\small
  \item Attention Is All You Need (2017). \texttt{arXiv:1706.03762}
  \item Proximal Policy Optimization Algorithms (2017). \texttt{arXiv:1707.06347}
  \item Language Models are Few-Shot Learners (2020). \texttt{arXiv:2005.14165}
  \item Scaling Laws for Neural Language Models (2020). \texttt{arXiv:2001.08361}
  \item Measuring Massive Multitask Language Understanding (2020). \texttt{arXiv:2009.03300}
  \item Learning to summarize from human feedback (2020). \texttt{arXiv:2009.01325}
  \item Evaluating Large Language Models Trained on Code (2021). \texttt{arXiv:2107.03374}
  \item Training Verifiers to Solve Math Word Problems (2021). \texttt{arXiv:2110.14168}
  \item Measuring Mathematical Problem Solving With the MATH Dataset (2021). \texttt{arXiv:2103.03874}
  \item Program Synthesis with Large Language Models (2021). \texttt{arXiv:2108.07732}
  \item Show Your Work: Scratchpads for Intermediate Computation with Language
  Models (2021). \texttt{arXiv:2112.00114}
  \item Chain-of-Thought Prompting Elicits Reasoning in Large Language Models (2022). \texttt{arXiv:2201.11903}
  \item Self-Consistency Improves Chain of Thought Reasoning in Language Models (2022). \texttt{arXiv:2203.11171}
  \item Training language models to follow instructions with human feedback (2022). \texttt{arXiv:2203.02155}
  \item Large Language Models are Zero-Shot Reasoners (2022). \texttt{arXiv:2205.11916}
  \item STaR: Bootstrapping Reasoning With Reasoning (2022). \texttt{arXiv:2203.14465}
  \item Solving Quantitative Reasoning Problems with Language Models (2022). \texttt{arXiv:2206.14858}
  \item Least-to-Most Prompting Enables Complex Reasoning in Large Language
  Models (2022). \texttt{arXiv:2205.10625}
  \item Solving math word problems with process- and outcome-based feedback (2022). \texttt{arXiv:2211.14275}
  \item Training Compute-Optimal Large Language Models (2022). \texttt{arXiv:2203.15556}
  \item Making Large Language Models Better Reasoners with Step-Aware Verifier (2022). \texttt{arXiv:2206.02336}
  \item Training a Helpful and Harmless Assistant with Reinforcement Learning
  from Human Feedback (2022). \texttt{arXiv:2204.05862}
  \item ReAct: Synergizing Reasoning and Acting in Language Models (2022). \texttt{arXiv:2210.03629}
  \item Competition-Level Code Generation with AlphaCode (2022). \texttt{arXiv:2203.07814}
  \item PAL: Program-aided Language Models (2022). \texttt{arXiv:2211.10435}
  \item GPT-4 Technical Report (2023). \texttt{arXiv:2303.08774}
  \item Tree of Thoughts: Deliberate Problem Solving with Large Language Models (2023). \texttt{arXiv:2305.10601}
  \item LLaMA: Open and Efficient Foundation Language Models (2023). \texttt{arXiv:2302.13971}
  \item Self-Refine: Iterative Refinement with Self-Feedback (2023). \texttt{arXiv:2303.17651}
  \item Llama 2: Open Foundation and Fine-Tuned Chat Models (2023). \texttt{arXiv:2307.09288}
  \item Let's Verify Step by Step (2023). \texttt{arXiv:2305.20050}
  \item Math-Shepherd: Verify and Reinforce LLMs Step-by-step without Human
  Annotations (2023). \texttt{arXiv:2312.08935}
  \item GPQA: A Graduate-Level Google-Proof Q\&A Benchmark (2023). \texttt{arXiv:2311.12022}
  \item Efficient Memory Management for Large Language Model Serving with
  PagedAttention (2023). \texttt{arXiv:2309.06180}
  \item Reflexion: Language Agents with Verbal Reinforcement Learning (2023). \texttt{arXiv:2303.11366}
  \item Alphazero-like Tree-Search can Guide Large Language Model Decoding and
  Training (2023). \texttt{arXiv:2309.17179}
  \item Direct Preference Optimization: Your Language Model is Secretly a Reward
  Model (2023). \texttt{arXiv:2305.18290}
  \item Reasoning with Language Model is Planning with World Model (2023). \texttt{arXiv:2305.14992}
  \item Scaling LLM Test-Time Compute Optimally can be More Effective than
  Scaling Model Parameters (2024). \texttt{arXiv:2408.03314}
  \item DeepSeekMath: Pushing the Limits of Mathematical Reasoning in Open
  Language Models (2024). \texttt{arXiv:2402.03300}
  \item The Llama 3 Herd of Models (2024). \texttt{arXiv:2407.21783}
  \item Large Language Monkeys: Scaling Inference Compute with Repeated Sampling (2024). \texttt{arXiv:2407.21787}
  \item OpenAI o1 System Card (2024). \texttt{arXiv:2412.16720}
  \item Qwen2.5 Technical Report (2024). \texttt{arXiv:2412.15115}
  \item Inference Scaling Laws: An Empirical Analysis of Compute-Optimal
  Inference for Problem-Solving with Language Models (2024). \texttt{arXiv:2408.00724}
  \item Do NOT Think That Much for 2+3=? On the Overthinking of o1-Like LLMs (2024). \texttt{arXiv:2412.21187}
  \item The Art of Scaling Test-Time Compute for Large Language Models (2025). \texttt{arXiv:2512.02008}
  \item Scaling over Scaling: Exploring Test-Time Scaling Plateau in Large Reasoning Models (2025). \texttt{arXiv:2505.20522}
  \item Scaling Test-time Compute for LLM Agents (2025). \texttt{arXiv:2506.12928}
  \item The Energy Cost of Reasoning: Analyzing Energy Usage in LLMs with Test-time Compute (2025). \texttt{arXiv:2505.14733}
  \item A Survey on Test-Time Scaling in Large Language Models: What, How,
  Where, and How Well? (2025). \texttt{arXiv:2503.24235}
  \item Are We Scaling the Right Thing? A System Perspective on Test-Time Scaling (2025). \texttt{arXiv:2509.19645}
  \item Rethinking Fine-Tuning when Scaling Test-Time Compute: Limiting Confidence Improves Mathematical Reasoning (2025). \texttt{arXiv:2502.07154}
  \item m1: Unleash the Potential of Test-Time Scaling for Medical Reasoning with Large Language Models (2025). \texttt{arXiv:2504.00869}
  \item Scaling up Test-Time Compute with Latent Reasoning: A Recurrent Depth
  Approach (2025). \texttt{arXiv:2502.05171}
  \item Rethinking Test-Time Scaling for Medical AI: Model and Task-Aware Strategies for LLMs and VLMs (2025). \texttt{arXiv:2506.13102}
  \item Latency and Token-Aware Test-Time Compute (2025). \texttt{arXiv:2509.09864}
  \item Inference-Time Scaling for Complex Tasks: Where We Stand and What Lies
  Ahead (2025). \texttt{arXiv:2504.00294}
  \item Reasoning on a Budget: A Survey of Adaptive and Controllable Test-Time Compute in LLMs (2025). \texttt{arXiv:2507.02076}
  \item Test-time Scaling of LLMs: A Survey from A Subproblem Structure Perspective (2025). \texttt{arXiv:2511.14772}
  \item Test-Time Scaling in Reasoning Models Is Not Effective for Knowledge-Intensive Tasks Yet (2025). \texttt{arXiv:2509.06861}
  \item Test-Time Scaling of Reasoning Models for Machine Translation (2025). \texttt{arXiv:2510.06471}
  \item Scaling Up, Speeding Up: A Benchmark of Speculative Decoding for Efficient LLM Test-Time Scaling (2025). \texttt{arXiv:2509.04474}
  \item Can 1B LLM Surpass 405B LLM? Rethinking Compute-Optimal Test-Time
  Scaling (2025). \texttt{arXiv:2502.06703}
  \item Inverse Scaling in Test-Time Compute (2025). \texttt{arXiv:2507.14417}
  \item Scaling Evaluation-time Compute with Reasoning Models as Evaluators (2025). \texttt{arXiv:2503.19877}
  \item Understanding the Role of Training Data in Test-Time Scaling (2025). \texttt{arXiv:2510.03605}
  \item Towards Thinking-Optimal Scaling of Test-Time Compute for LLM Reasoning (2025). \texttt{arXiv:2502.18080}
  \item On the Role of Temperature Sampling in Test-Time Scaling (2025). \texttt{arXiv:2510.02611}
  \item Prompting Test-Time Scaling Is A Strong LLM Reasoning Data Augmentation (2025). \texttt{arXiv:2510.09599}
  \item Inference-Time Computations for LLM Reasoning and Planning: A Benchmark
  and Insights (2025). \texttt{arXiv:2502.12521}
  \item Limits and Gains of Test-Time Scaling in Vision-Language Reasoning (2025). \texttt{arXiv:2512.11109}
  \item Adaptive Test-Time Reasoning via Reward-Guided Dual-Phase Search (2025). \texttt{arXiv:2509.25420}
  \item ScaleRTL: Scaling LLMs with Reasoning Data and Test-Time Compute for Accurate RTL Code Generation (2025). \texttt{arXiv:2506.05566}
  \item AgentTTS: Large Language Model Agent for Test-time Compute-optimal Scaling Strategy in Complex Tasks (2025). \texttt{arXiv:2508.00890}
  \item Efficient Test-Time Scaling via Self-Calibration (2025). \texttt{arXiv:2503.00031}
  \item Enhancing Test-Time Scaling of Large Language Models with Hierarchical Retrieval-Augmented MCTS (2025). \texttt{arXiv:2507.05557}
  \item DeepSeek-R1: Incentivizing Reasoning Capability in LLMs via Reinforcement Learning (2025). \texttt{arXiv:2501.12948}
  \item s1: Simple test-time scaling (2025). \texttt{arXiv:2501.19393}
  \item TMAS: Scaling Test-Time Compute via Multi-Agent Synergy (2026). \texttt{arXiv:2605.10344}
\end{enumerate}

\end{document}
